\documentclass[11pt,a4paper]{article}
\usepackage[utf8]{inputenc}
\usepackage[T1]{fontenc}
\usepackage[spanish,es-minimal,es-noshorthands]{babel}
\usepackage{lmodern}
\usepackage[sfdefault,lining]{FiraSans}
\usepackage[a4paper,top=2.3cm,bottom=2.2cm,left=1.8cm,right=1.8cm,headheight=1.3cm,headsep=0.4cm,footskip=1cm]{geometry}
\usepackage{amsmath,amssymb}
\usepackage{icomma}
\usepackage{graphicx}
\usepackage[table]{xcolor}
\usepackage{tikz}
\usepackage{fontawesome5}
\usepackage{ccicons}
\usepackage{fancyhdr}
\usepackage{lastpage}
\usepackage{tasks}
\usepackage{needspace}
\usepackage[most]{tcolorbox}
\usepackage{array}
\usepackage[hidelinks]{hyperref}
\AtBeginDocument{\hypersetup{pdftitle={Ficha de ejercicios de la unidad 11: Cuerpos geométricos. 3º ESO},pdfsubject={Matemáticas 3º ESO},pdfkeywords={poliedros, fórmula de Euler, poliedros regulares, prisma, pirámide, cilindro, cono, esfera, área, volumen, capacidad}}}

\definecolor{navy}{HTML}{1B2748}
\definecolor{accent}{HTML}{E8603A}
\definecolor{teal}{HTML}{1F9E90}
\definecolor{soft}{HTML}{EEF1F8}
\definecolor{muted}{HTML}{5D6685}

\setlength{\parindent}{0pt}
\setlength{\parskip}{3pt}
\settasks{label=\alph*), label-format=\bfseries\color{navy}, label-width=1.3em, item-indent=1.8em, column-sep=0.6em, after-item-skip=5pt, before-skip=3pt, after-skip=4pt}

% ---------- capçalera i peu ----------
\newcommand{\llicencia}{\href{https://creativecommons.org/licenses/by-nc-sa/4.0/deed.es}{\ccbyncsa}\ \ Licencia Creative Commons Atribución-NoComercial-CompartirIgual 4.0 Internacional (CC BY-NC-SA 4.0)}
\pagestyle{fancy}
\fancyhf{}
\fancyhead[L]{\raisebox{-0.35\height}{\includegraphics[height=0.95cm]{logo-gva.png}}}
\fancyhead[C]{\small\color{muted}Ficha de ejercicios · Unidad 11: Cuerpos geométricos · 3\textsuperscript{o} ESO}
\fancyhead[R]{\raisebox{-0.35\height}{\includegraphics[height=0.95cm]{logo-vermellar.png}}}
\fancyfoot[L]{\scriptsize\color{muted}\llicencia}
\fancyfoot[R]{\small\color{muted}\thepage}
\renewcommand{\headrulewidth}{0.4pt}
\renewcommand{\headrule}{\color{muted!40}\hrule height \headrulewidth}
\fancypagestyle{primera}{\fancyhf{}\renewcommand{\headrulewidth}{0pt}
  \fancyfoot[L]{\scriptsize\color{muted}\llicencia}
  \fancyfoot[R]{\small\color{muted}\thepage}}

% ---------- nivells i exercicis ----------
\newcommand{\estrelles}[1]{{\color{accent}\ifcase#1\or\faStar\faStar[regular]\faStar[regular]\or\faStar\faStar\faStar[regular]\or\faStar\faStar\faStar\fi}}
\newcounter{exer}
\newcommand{\ex}[1]{\par\needspace{4\baselineskip}\medskip\stepcounter{exer}\noindent{\color{navy}\bfseries\theexer.}\ \ {\small\estrelles{#1}}\ \ }

\newcommand{\seccio}[2]{\par\needspace{10\baselineskip}\bigskip
  \begin{tcolorbox}[colback=navy,colframe=navy,coltext=white,arc=2mm,boxrule=0pt,left=3mm,right=3mm,top=1.2mm,bottom=1.2mm]
  {\large\bfseries #1\quad #2}
  \end{tcolorbox}\vspace{-2pt}}

\newcolumntype{C}[1]{>{\centering\arraybackslash}p{#1}}
\usepackage{pifont}

% quadre de recordatori (només a la fitxa d'exercicis)
\newcommand{\recorda}[1]{\par\smallskip\begin{tcolorbox}[colback=accent!7,colframe=accent!55,arc=2mm,boxrule=0.5pt,left=2.5mm,right=2.5mm,top=1mm,bottom=1mm]\small{\color{accent}\faLightbulb[regular]}\ \textbf{Recuerda:}\ #1\end{tcolorbox}}
\newcommand{\res}[1]{#1}
\newcommand{\rest}[1]{#1}

\def\fmneed{14}
% =====================================================================
% Capçalera FentMates de les fitxes: un color per a cada curs.
% Els colors són els de la pàgina inicial (LEVELS a assets/js/materials.js);
% el logo ve de _eines/logo/pdf/ (python3 _eines/logo/genera.py).
%
% Ús, al final del preàmbul de la unitat:
%   % =====================================================================
% Capçalera FentMates de les fitxes: un color per a cada curs.
% Els colors són els de la pàgina inicial (LEVELS a assets/js/materials.js);
% el logo ve de _eines/logo/pdf/ (python3 _eines/logo/genera.py).
%
% Ús, al final del preàmbul de la unitat:
%   % =====================================================================
% Capçalera FentMates de les fitxes: un color per a cada curs.
% Els colors són els de la pàgina inicial (LEVELS a assets/js/materials.js);
% el logo ve de _eines/logo/pdf/ (python3 _eines/logo/genera.py).
%
% Ús, al final del preàmbul de la unitat:
%   \input{../comu/fentmates.tex}
%   \fmcurs{1eso}
% i, al començament del document:
%   \capcalera{Matemàtiques}{1r ESO}{Fitxa d'exercicis · Unitat 2}{Divisibilitat}
% Opcional abans de \input: \def\fmneed{14} (espai mínim abans de cada secció).
% =====================================================================
\makeatletter
\def\fm@logodir{../../logo/pdf}
\providecommand{\fmneed}{10}
\def\fm@cursnom{}

% clau del curs, color i color del text damunt del color
\newcommand{\fm@curs}[3]{\@namedef{fm@color@#1}{#2}\@namedef{fm@text@#1}{#3}}
\fm@curs{1eso}{E6B84A}{navy}
\fm@curs{2eso}{3B82F6}{white}
\fm@curs{3eso}{FF7A50}{white}
\fm@curs{4eso-a}{A78BFA}{white}
\fm@curs{4eso-b}{6D28D9}{white}
\fm@curs{1batx-mi}{9B1C1C}{white}
\fm@curs{1batx-ccss}{EF6B6B}{white}
\fm@curs{2batx-mii}{1E6B47}{white}
\fm@curs{2batx-ccss}{2BB3A3}{white}

\newcommand{\fmcurs}[1]{%
  \@ifundefined{fm@color@#1}{\PackageError{fentmates}{Curs desconegut: #1}{}}{}%
  \edef\fm@tmp{\noexpand\definecolor{curs}{HTML}{\@nameuse{fm@color@#1}}%
               \noexpand\colorlet{curstext}{\@nameuse{fm@text@#1}}}\fm@tmp
  \colorlet{cursfosc}{curs!72!black}%
  \def\fm@logo{\fm@logodir/vertical-#1.pdf}%
  \def\fm@marca{\fm@logodir/marca-#1.pdf}}

% ---------- capçalera de la primera pàgina ----------
% #1 matèria (petit), #2 curs, #3 tipus de document i unitat, #4 títol
\newcommand{\capcalera}[4]{%
  \gdef\fm@cursnom{#2}%
  \par\noindent
  \begin{tcolorbox}[colback=curs,colframe=curs,boxrule=0pt,arc=2.2mm,width=3.5cm,height=2.15cm,
      valign=center,left=3mm,right=2mm,top=1mm,bottom=1mm,nobeforeafter,box align=center]
    \color{curstext}{\footnotesize #1}\par\vspace{1pt}{\fontsize{19}{21}\selectfont\bfseries #2\par}%
  \end{tcolorbox}\hspace{5mm}%
  \begin{minipage}[c]{\dimexpr\linewidth-3.5cm-5mm-2.9cm\relax}
    {\color{cursfosc}\small\bfseries\MakeUppercase{#3}}\par\vspace{2pt}
    {\color{navy}\fontsize{20}{23}\selectfont\bfseries #4\par}
  \end{minipage}\hfill
  \begin{minipage}[c]{2.6cm}\raggedleft\includegraphics[width=2.6cm]{\fm@logo}\end{minipage}\par
  \vspace{2pt}}

% ---------- barres de secció del color del curs ----------
\renewcommand{\seccio}[2]{\par\needspace{\fmneed\baselineskip}\bigskip
  \begin{tcolorbox}[colback=curs,colframe=curs,coltext=curstext,arc=2mm,boxrule=0pt,left=3mm,right=3mm,top=1.2mm,bottom=1.2mm]
  {\large\bfseries #1\quad #2}
  \end{tcolorbox}\vspace{-2pt}}

% ---------- estrelles de dificultat del color del curs (a Batxillerat no se'n posen) ----------
\@ifundefined{estrelles}{}{%
  \renewcommand{\estrelles}[1]{{\color{curs}\ifcase#1\or\faStar\faStar[regular]\faStar[regular]\or\faStar\faStar\faStar[regular]\or\faStar\faStar\faStar\fi}}}

% ---------- capçalera i peu de les pàgines ----------
% primera pàgina: logos institucionals al peu
\fancypagestyle{primera}{\fancyhf{}\renewcommand{\headrulewidth}{0pt}
  \fancyfoot[L]{\raisebox{-0.32\height}{\includegraphics[height=0.78cm]{logo-gva.png}}\hspace{5mm}%
    \raisebox{-0.32\height}{\includegraphics[height=0.78cm]{logo-vermellar.png}}}
  \fancyfoot[R]{\scriptsize\color{muted}\href{https://creativecommons.org/licenses/by-nc-sa/4.0/deed.ca}{\ccbyncsa}\ \ CC BY-NC-SA 4.0 · fentmates.org\qquad\small\thepage}}
% la resta: el símbol a l'esquerra i el curs a la dreta, amb el títol al centre com abans
\fancyhead[L]{\raisebox{-0.3\height}{\includegraphics[height=0.62cm]{\fm@marca}}}
\fancyhead[R]{\small\bfseries\color{cursfosc}\fm@cursnom}
\renewcommand{\headrule}{\color{curs!45}\hrule height \headrulewidth}
\makeatother

%   \fmcurs{1eso}
% i, al començament del document:
%   \capcalera{Matemàtiques}{1r ESO}{Fitxa d'exercicis · Unitat 2}{Divisibilitat}
% Opcional abans de \input: \def\fmneed{14} (espai mínim abans de cada secció).
% =====================================================================
\makeatletter
\def\fm@logodir{../../logo/pdf}
\providecommand{\fmneed}{10}
\def\fm@cursnom{}

% clau del curs, color i color del text damunt del color
\newcommand{\fm@curs}[3]{\@namedef{fm@color@#1}{#2}\@namedef{fm@text@#1}{#3}}
\fm@curs{1eso}{E6B84A}{navy}
\fm@curs{2eso}{3B82F6}{white}
\fm@curs{3eso}{FF7A50}{white}
\fm@curs{4eso-a}{A78BFA}{white}
\fm@curs{4eso-b}{6D28D9}{white}
\fm@curs{1batx-mi}{9B1C1C}{white}
\fm@curs{1batx-ccss}{EF6B6B}{white}
\fm@curs{2batx-mii}{1E6B47}{white}
\fm@curs{2batx-ccss}{2BB3A3}{white}

\newcommand{\fmcurs}[1]{%
  \@ifundefined{fm@color@#1}{\PackageError{fentmates}{Curs desconegut: #1}{}}{}%
  \edef\fm@tmp{\noexpand\definecolor{curs}{HTML}{\@nameuse{fm@color@#1}}%
               \noexpand\colorlet{curstext}{\@nameuse{fm@text@#1}}}\fm@tmp
  \colorlet{cursfosc}{curs!72!black}%
  \def\fm@logo{\fm@logodir/vertical-#1.pdf}%
  \def\fm@marca{\fm@logodir/marca-#1.pdf}}

% ---------- capçalera de la primera pàgina ----------
% #1 matèria (petit), #2 curs, #3 tipus de document i unitat, #4 títol
\newcommand{\capcalera}[4]{%
  \gdef\fm@cursnom{#2}%
  \par\noindent
  \begin{tcolorbox}[colback=curs,colframe=curs,boxrule=0pt,arc=2.2mm,width=3.5cm,height=2.15cm,
      valign=center,left=3mm,right=2mm,top=1mm,bottom=1mm,nobeforeafter,box align=center]
    \color{curstext}{\footnotesize #1}\par\vspace{1pt}{\fontsize{19}{21}\selectfont\bfseries #2\par}%
  \end{tcolorbox}\hspace{5mm}%
  \begin{minipage}[c]{\dimexpr\linewidth-3.5cm-5mm-2.9cm\relax}
    {\color{cursfosc}\small\bfseries\MakeUppercase{#3}}\par\vspace{2pt}
    {\color{navy}\fontsize{20}{23}\selectfont\bfseries #4\par}
  \end{minipage}\hfill
  \begin{minipage}[c]{2.6cm}\raggedleft\includegraphics[width=2.6cm]{\fm@logo}\end{minipage}\par
  \vspace{2pt}}

% ---------- barres de secció del color del curs ----------
\renewcommand{\seccio}[2]{\par\needspace{\fmneed\baselineskip}\bigskip
  \begin{tcolorbox}[colback=curs,colframe=curs,coltext=curstext,arc=2mm,boxrule=0pt,left=3mm,right=3mm,top=1.2mm,bottom=1.2mm]
  {\large\bfseries #1\quad #2}
  \end{tcolorbox}\vspace{-2pt}}

% ---------- estrelles de dificultat del color del curs (a Batxillerat no se'n posen) ----------
\@ifundefined{estrelles}{}{%
  \renewcommand{\estrelles}[1]{{\color{curs}\ifcase#1\or\faStar\faStar[regular]\faStar[regular]\or\faStar\faStar\faStar[regular]\or\faStar\faStar\faStar\fi}}}

% ---------- capçalera i peu de les pàgines ----------
% primera pàgina: logos institucionals al peu
\fancypagestyle{primera}{\fancyhf{}\renewcommand{\headrulewidth}{0pt}
  \fancyfoot[L]{\raisebox{-0.32\height}{\includegraphics[height=0.78cm]{logo-gva.png}}\hspace{5mm}%
    \raisebox{-0.32\height}{\includegraphics[height=0.78cm]{logo-vermellar.png}}}
  \fancyfoot[R]{\scriptsize\color{muted}\href{https://creativecommons.org/licenses/by-nc-sa/4.0/deed.ca}{\ccbyncsa}\ \ CC BY-NC-SA 4.0 · fentmates.org\qquad\small\thepage}}
% la resta: el símbol a l'esquerra i el curs a la dreta, amb el títol al centre com abans
\fancyhead[L]{\raisebox{-0.3\height}{\includegraphics[height=0.62cm]{\fm@marca}}}
\fancyhead[R]{\small\bfseries\color{cursfosc}\fm@cursnom}
\renewcommand{\headrule}{\color{curs!45}\hrule height \headrulewidth}
\makeatother

%   \fmcurs{1eso}
% i, al començament del document:
%   \capcalera{Matemàtiques}{1r ESO}{Fitxa d'exercicis · Unitat 2}{Divisibilitat}
% Opcional abans de \input: \def\fmneed{14} (espai mínim abans de cada secció).
% =====================================================================
\makeatletter
\def\fm@logodir{../../logo/pdf}
\providecommand{\fmneed}{10}
\def\fm@cursnom{}

% clau del curs, color i color del text damunt del color
\newcommand{\fm@curs}[3]{\@namedef{fm@color@#1}{#2}\@namedef{fm@text@#1}{#3}}
\fm@curs{1eso}{E6B84A}{navy}
\fm@curs{2eso}{3B82F6}{white}
\fm@curs{3eso}{FF7A50}{white}
\fm@curs{4eso-a}{A78BFA}{white}
\fm@curs{4eso-b}{6D28D9}{white}
\fm@curs{1batx-mi}{9B1C1C}{white}
\fm@curs{1batx-ccss}{EF6B6B}{white}
\fm@curs{2batx-mii}{1E6B47}{white}
\fm@curs{2batx-ccss}{2BB3A3}{white}

\newcommand{\fmcurs}[1]{%
  \@ifundefined{fm@color@#1}{\PackageError{fentmates}{Curs desconegut: #1}{}}{}%
  \edef\fm@tmp{\noexpand\definecolor{curs}{HTML}{\@nameuse{fm@color@#1}}%
               \noexpand\colorlet{curstext}{\@nameuse{fm@text@#1}}}\fm@tmp
  \colorlet{cursfosc}{curs!72!black}%
  \def\fm@logo{\fm@logodir/vertical-#1.pdf}%
  \def\fm@marca{\fm@logodir/marca-#1.pdf}}

% ---------- capçalera de la primera pàgina ----------
% #1 matèria (petit), #2 curs, #3 tipus de document i unitat, #4 títol
\newcommand{\capcalera}[4]{%
  \gdef\fm@cursnom{#2}%
  \par\noindent
  \begin{tcolorbox}[colback=curs,colframe=curs,boxrule=0pt,arc=2.2mm,width=3.5cm,height=2.15cm,
      valign=center,left=3mm,right=2mm,top=1mm,bottom=1mm,nobeforeafter,box align=center]
    \color{curstext}{\footnotesize #1}\par\vspace{1pt}{\fontsize{19}{21}\selectfont\bfseries #2\par}%
  \end{tcolorbox}\hspace{5mm}%
  \begin{minipage}[c]{\dimexpr\linewidth-3.5cm-5mm-2.9cm\relax}
    {\color{cursfosc}\small\bfseries\MakeUppercase{#3}}\par\vspace{2pt}
    {\color{navy}\fontsize{20}{23}\selectfont\bfseries #4\par}
  \end{minipage}\hfill
  \begin{minipage}[c]{2.6cm}\raggedleft\includegraphics[width=2.6cm]{\fm@logo}\end{minipage}\par
  \vspace{2pt}}

% ---------- barres de secció del color del curs ----------
\renewcommand{\seccio}[2]{\par\needspace{\fmneed\baselineskip}\bigskip
  \begin{tcolorbox}[colback=curs,colframe=curs,coltext=curstext,arc=2mm,boxrule=0pt,left=3mm,right=3mm,top=1.2mm,bottom=1.2mm]
  {\large\bfseries #1\quad #2}
  \end{tcolorbox}\vspace{-2pt}}

% ---------- estrelles de dificultat del color del curs (a Batxillerat no se'n posen) ----------
\@ifundefined{estrelles}{}{%
  \renewcommand{\estrelles}[1]{{\color{curs}\ifcase#1\or\faStar\faStar[regular]\faStar[regular]\or\faStar\faStar\faStar[regular]\or\faStar\faStar\faStar\fi}}}

% ---------- capçalera i peu de les pàgines ----------
% primera pàgina: logos institucionals al peu
\fancypagestyle{primera}{\fancyhf{}\renewcommand{\headrulewidth}{0pt}
  \fancyfoot[L]{\raisebox{-0.32\height}{\includegraphics[height=0.78cm]{logo-gva.png}}\hspace{5mm}%
    \raisebox{-0.32\height}{\includegraphics[height=0.78cm]{logo-vermellar.png}}}
  \fancyfoot[R]{\scriptsize\color{muted}\href{https://creativecommons.org/licenses/by-nc-sa/4.0/deed.ca}{\ccbyncsa}\ \ CC BY-NC-SA 4.0 · fentmates.org\qquad\small\thepage}}
% la resta: el símbol a l'esquerra i el curs a la dreta, amb el títol al centre com abans
\fancyhead[L]{\raisebox{-0.3\height}{\includegraphics[height=0.62cm]{\fm@marca}}}
\fancyhead[R]{\small\bfseries\color{cursfosc}\fm@cursnom}
\renewcommand{\headrule}{\color{curs!45}\hrule height \headrulewidth}
\makeatother

\fmcurs{3eso}

\begin{document}
\thispagestyle{primera}

\capcalera{Matemáticas}{3\textsuperscript{o} ESO}{Ficha de ejercicios · Unidad 11}{Cuerpos geométricos}
\vspace{10pt}
Nombre y apellidos: \rule{6.6cm}{0.4pt}\hfill Grupo: \rule{1.6cm}{0.4pt}\hfill Fecha: \rule{2.2cm}{0.4pt}
\vspace{8pt}

\begin{tcolorbox}[colback=soft,colframe=soft,arc=2mm,left=3mm,right=3mm,top=1.5mm,bottom=1.5mm]
\small Ejercicios graduados por niveles de dificultad:\quad \estrelles{1}\ básico \qquad \estrelles{2}\ medio \qquad \estrelles{3}\ avanzado.\\[2pt] Resuélvelos en el cuaderno mostrando todos los pasos y haz un dibujo de cada cuerpo. Si un resultado no es exacto, redondéalo a las centésimas; para $\pi$, usa la tecla de la calculadora.
\end{tcolorbox}

% =====================================================
\seccio{1}{Poliedros. Fórmula de Euler. Poliedros regulares}

\ex{1} Copia y completa la tabla. Comprueba en cada caso la fórmula de Euler.\par\smallskip\begin{center}\renewcommand{\arraystretch}{1.25}\begin{tabular}{|l|C{1.6cm}|C{1.6cm}|C{1.6cm}|}\hline Poliedro & Caras & Vértices & Aristas \\\hline prisma triangular & & &  \\\hline pirámide pentagonal & & &  \\\hline prisma hexagonal & & &  \\\hline pirámide octogonal & & &  \\\hline cubo & & &  \\\hline\end{tabular}\end{center}\recorda{En un poliedro convexo, \textbf{caras} + \textbf{vértices} = \textbf{aristas} + 2: \ $C+V=A+2$ (fórmula de Euler). Un prisma de base de $n$ lados tiene $n+2$ caras, $2n$ vértices y $3n$ aristas; una pirámide, $n+1$ caras, $n+1$ vértices y $2n$ aristas.}

\ex{1} Usa la fórmula de Euler para calcular lo que falta en cada poliedro convexo:\begin{tasks}(3)
  \task $C=8$, \ $V=12$, \ $A=\ ?$
  \task $V=10$, \ $A=15$, \ $C=\ ?$
  \task $C=12$, \ $A=30$, \ $V=\ ?$
\end{tasks}

\ex{1} ¿Cuáles de estos poliedros son regulares? Justifícalo.\begin{tasks}(3)
  \task tetraedro
  \task ortoedro
  \task octaedro
  \task pirámide cuadrangular regular
  \task icosaedro
  \task cubo
\end{tasks}

\ex{2} Un prisma tiene 24 aristas. ¿Qué polígono tiene por base? ¿Cuántas caras y cuántos vértices tiene?

\ex{2} Una pirámide tiene 11 vértices. ¿Qué polígono tiene por base? ¿Cuántas caras y aristas tiene?

\ex{2} ¿Puede existir un poliedro convexo con: \ a) 7 caras, 7 vértices y 12 aristas? \ b) 5 caras, 6 vértices y 10 aristas?

\ex{3} Si juntamos dos pirámides cuadrangulares iguales por la base obtenemos una \emph{bipirámide}. ¿Cuántas caras, vértices y aristas tiene? ¿Y si las pirámides son pentagonales?

\ex{3} Explica por qué no puede haber ningún poliedro regular con caras hexagonales.


% =====================================================
\seccio{2}{Prismas: elementos y área}

\ex{1} \begin{minipage}[t]{0.52\linewidth}Calcula el área total:\begin{tasks}(1)
  \task de un cubo de 4 cm de arista;
  \task del ortoedro de la figura.
\end{tasks}\end{minipage}\hfill\begin{minipage}[t]{0.34\linewidth}\vspace{0pt}\centering \begin{tikzpicture}[x=0.45cm,y=0.45cm,baseline=(current bounding box.center),line join=round]
\path[fill=teal!14,draw=none] (0.000,3.709) -- (0.000,0.000) -- (5.298,-1.055) -- (5.298,2.654) -- cycle;
\path[fill=teal!14,draw=none] (5.298,2.654) -- (5.298,-1.055) -- (6.706,-0.063) -- (6.706,3.646) -- cycle;
\path[fill=teal!14,draw=none] (1.408,4.701) -- (0.000,3.709) -- (5.298,2.654) -- (6.706,3.646) -- cycle;
\draw[dashed,draw=muted] (0.000,0.000) -- (1.408,0.992);
\draw[dashed,draw=muted] (1.408,0.992) -- (6.706,-0.063);
\draw[dashed,draw=muted] (1.408,0.992) -- (1.408,4.701);
\draw[thick] (0.000,0.000) -- (5.298,-1.055);
\draw[thick] (5.298,-1.055) -- (6.706,-0.063);
\draw[thick] (0.000,0.000) -- (0.000,3.709);
\draw[thick] (5.298,-1.055) -- (5.298,2.654);
\draw[thick] (0.000,3.709) -- (5.298,2.654);
\draw[thick] (0.000,3.709) -- (1.408,4.701);
\draw[thick] (6.706,-0.063) -- (6.706,3.646);
\draw[thick] (5.298,2.654) -- (6.706,3.646);
\draw[thick] (1.408,4.701) -- (6.706,3.646);
\node[font=\small,inner sep=2pt,below] at (2.649,-0.528) {$6$\,cm};
\node[font=\small,inner sep=2pt,right] at (6.002,-0.559) {$3$\,cm};
\node[font=\small,inner sep=2pt,right] at (5.298,0.799) {$4$\,cm};
\end{tikzpicture}\end{minipage}\par\recorda{Área de un prisma recto: $A_L=P_{\text{base}}\cdot h$, \ $A_T=A_L+2\,A_{\text{base}}$. Ortoedro de aristas $a,\,b,\,c$: $A_T=2(ab+ac+bc)$ y la diagonal es $D=\sqrt{a^2+b^2+c^2}$.}

\ex{1} Calcula la diagonal de un ortoedro de 12 cm × 4 cm × 3 cm.

\ex{1} Un prisma cuadrangular regular tiene 5 cm de arista de la base y 12 cm de altura. Calcula su área lateral y su área total.

\ex{2} Calcula el área total de un prisma triangular regular de 4 cm de arista de la base y 10 cm de altura.

\ex{2} Calcula el área total de un prisma hexagonal regular de 6 cm de arista de la base y 8 cm de altura.

\ex{2} Queremos hacer una caja sin tapa de 40 cm × 30 cm de base y 20 cm de altura. ¿Cuánto cartón necesitamos?

\ex{3} La diagonal de un cubo mide $\sqrt{48}$ cm. ¿Cuánto mide la arista? Calcula su área total.

\ex{3} La base de un prisma recto es un triángulo rectángulo de catetos 6 cm y 8 cm, y la altura del prisma es 15 cm. Calcula su área total.


% =====================================================
\seccio{3}{Pirámides y troncos de pirámide}

\ex{1} \begin{minipage}[t]{0.56\linewidth}La pirámide cuadrangular regular de la figura tiene 10 cm de arista de la base y 12 cm de altura. Calcula su apotema y su área total.\end{minipage}\hfill\begin{minipage}[t]{0.3\linewidth}\vspace{0pt}\centering \begin{tikzpicture}[x=0.5cm,y=0.5cm,baseline=(current bounding box.center),line join=round]
\path[fill=teal!14,draw=none] (0.000,3.338) -- (-2.029,-0.232) -- (0.620,-0.760) -- cycle;
\path[fill=teal!14,draw=none] (0.000,3.338) -- (0.620,-0.760) -- (2.029,0.232) -- cycle;
\draw[dashed,draw=muted] (-2.029,-0.232) -- (-0.620,0.760);
\draw[dashed,draw=muted] (2.029,0.232) -- (-0.620,0.760);
\draw[dashed,draw=muted] (-0.620,0.760) -- (0.000,3.338);
\draw[thick] (-2.029,-0.232) -- (0.620,-0.760);
\draw[thick] (0.620,-0.760) -- (2.029,0.232);
\draw[thick] (-2.029,-0.232) -- (0.000,3.338);
\draw[thick] (0.620,-0.760) -- (0.000,3.338);
\draw[thick] (2.029,0.232) -- (0.000,3.338);
\node[font=\small,inner sep=2pt,below] at (-0.704,-0.496) {$10$\,cm};
\draw[dashed,draw=accent,thick] (0.000,0.000) -- (0.000,3.338);
\draw[dashed,draw=accent,thick] (0.000,0.000) -- (1.324,-0.264);
\draw[draw=navy,thick] (1.324,-0.264) -- (0.000,3.338);
\fill (0.000,0.000) circle (1.3pt);
\node[font=\small,inner sep=2pt,left,text=accent] at (0.000,1.669) {$12$\,cm};
\end{tikzpicture}\end{minipage}\par\recorda{En una pirámide regular, la altura $h$, la apotema de la base $a_b$ y la apotema de la pirámide $a_p$ forman un triángulo rectángulo: $a_p^2=h^2+a_b^2$. \ $A_L=\frac{P_{\text{base}}\cdot a_p}{2}$, \ $A_T=A_L+A_{\text{base}}$.}

\ex{1} Una pirámide hexagonal regular tiene 4 cm de arista de la base y 9 cm de apotema. Calcula su área lateral.

\ex{2} Una pirámide cuadrangular regular tiene 6 cm de arista de la base y 5 cm de apotema. ¿Qué altura tiene? Calcula su área total.

\ex{2} Una pirámide cuadrangular regular tiene 8 cm de arista de la base y 6 cm de altura. Calcula la apotema, la arista lateral y el área total.

\ex{2} Calcula el área total de un tetraedro regular de 6 cm de arista.

\ex{3} Un tronco de pirámide cuadrangular regular tiene las aristas de las bases de 12 cm y 6 cm y una altura de 4 cm. Calcula la apotema del tronco y el área total.

\ex{3} La gran pirámide de Keops (Guiza) tiene una base cuadrada de 230 m de lado y mide 146 m de altura. Calcula la apotema y el área de las cuatro caras laterales.


% =====================================================
\seccio{4}{Cilindro y cono}

\ex{1} \begin{minipage}[t]{0.61\linewidth}Calcula el área lateral y el área total del cilindro de la figura.\end{minipage}\hfill\begin{minipage}[t]{0.25\linewidth}\vspace{0pt}\centering \begin{tikzpicture}[x=0.5cm,y=0.5cm,baseline=(current bounding box.center),line join=round]
\path[fill=teal!14,draw=none] (-1.500,2.967) -- (-1.500,0.000) -- (-1.497,-0.037) -- (-1.487,-0.073) -- (-1.471,-0.110) -- (-1.449,-0.145) -- (-1.420,-0.181) -- (-1.386,-0.215) -- (-1.345,-0.249) -- (-1.299,-0.281) -- (-1.247,-0.312) -- (-1.190,-0.342) -- (-1.128,-0.370) -- (-1.061,-0.397) -- (-0.989,-0.422) -- (-0.913,-0.446) -- (-0.833,-0.467) -- (-0.750,-0.487) -- (-0.663,-0.504) -- (-0.574,-0.519) -- (-0.482,-0.532) -- (-0.388,-0.543) -- (-0.293,-0.551) -- (-0.196,-0.557) -- (-0.098,-0.561) -- (-0.000,-0.562) -- (0.098,-0.561) -- (0.196,-0.557) -- (0.293,-0.551) -- (0.388,-0.543) -- (0.482,-0.532) -- (0.574,-0.519) -- (0.663,-0.504) -- (0.750,-0.487) -- (0.833,-0.467) -- (0.913,-0.446) -- (0.989,-0.422) -- (1.061,-0.397) -- (1.128,-0.370) -- (1.190,-0.342) -- (1.247,-0.312) -- (1.299,-0.281) -- (1.345,-0.249) -- (1.386,-0.215) -- (1.420,-0.181) -- (1.449,-0.145) -- (1.471,-0.110) -- (1.487,-0.073) -- (1.497,-0.037) -- (1.500,-0.000) -- (1.500,2.967) -- (1.500,2.967) -- (1.497,2.930) -- (1.487,2.894) -- (1.471,2.857) -- (1.449,2.822) -- (1.420,2.786) -- (1.386,2.752) -- (1.345,2.718) -- (1.299,2.686) -- (1.247,2.655) -- (1.190,2.625) -- (1.128,2.596) -- (1.061,2.570) -- (0.989,2.545) -- (0.913,2.521) -- (0.833,2.500) -- (0.750,2.480) -- (0.663,2.463) -- (0.574,2.448) -- (0.482,2.435) -- (0.388,2.424) -- (0.293,2.416) -- (0.196,2.410) -- (0.098,2.406) -- (-0.000,2.405) -- (-0.098,2.406) -- (-0.196,2.410) -- (-0.293,2.416) -- (-0.388,2.424) -- (-0.482,2.435) -- (-0.574,2.448) -- (-0.663,2.463) -- (-0.750,2.480) -- (-0.833,2.500) -- (-0.913,2.521) -- (-0.989,2.545) -- (-1.061,2.570) -- (-1.128,2.596) -- (-1.190,2.625) -- (-1.247,2.655) -- (-1.299,2.686) -- (-1.345,2.718) -- (-1.386,2.752) -- (-1.420,2.786) -- (-1.449,2.822) -- (-1.471,2.857) -- (-1.487,2.894) -- (-1.497,2.930) -- (-1.500,2.967) -- cycle;
\draw[dashed,draw=muted] (1.500,0.000) -- (1.497,0.037) -- (1.487,0.073) -- (1.471,0.110) -- (1.449,0.145) -- (1.420,0.181) -- (1.386,0.215) -- (1.345,0.249) -- (1.299,0.281) -- (1.247,0.312) -- (1.190,0.342) -- (1.128,0.370) -- (1.061,0.397) -- (0.989,0.422) -- (0.913,0.446) -- (0.833,0.467) -- (0.750,0.487) -- (0.663,0.504) -- (0.574,0.519) -- (0.482,0.532) -- (0.388,0.543) -- (0.293,0.551) -- (0.196,0.557) -- (0.098,0.561) -- (0.000,0.562) -- (-0.098,0.561) -- (-0.196,0.557) -- (-0.293,0.551) -- (-0.388,0.543) -- (-0.482,0.532) -- (-0.574,0.519) -- (-0.663,0.504) -- (-0.750,0.487) -- (-0.833,0.467) -- (-0.913,0.446) -- (-0.989,0.422) -- (-1.061,0.397) -- (-1.128,0.370) -- (-1.190,0.342) -- (-1.247,0.312) -- (-1.299,0.281) -- (-1.345,0.249) -- (-1.386,0.215) -- (-1.420,0.181) -- (-1.449,0.145) -- (-1.471,0.110) -- (-1.487,0.073) -- (-1.497,0.037) -- (-1.500,0.000);
\draw[thick] (-1.500,0.000) -- (-1.497,-0.037) -- (-1.487,-0.073) -- (-1.471,-0.110) -- (-1.449,-0.145) -- (-1.420,-0.181) -- (-1.386,-0.215) -- (-1.345,-0.249) -- (-1.299,-0.281) -- (-1.247,-0.312) -- (-1.190,-0.342) -- (-1.128,-0.370) -- (-1.061,-0.397) -- (-0.989,-0.422) -- (-0.913,-0.446) -- (-0.833,-0.467) -- (-0.750,-0.487) -- (-0.663,-0.504) -- (-0.574,-0.519) -- (-0.482,-0.532) -- (-0.388,-0.543) -- (-0.293,-0.551) -- (-0.196,-0.557) -- (-0.098,-0.561) -- (-0.000,-0.562) -- (0.098,-0.561) -- (0.196,-0.557) -- (0.293,-0.551) -- (0.388,-0.543) -- (0.482,-0.532) -- (0.574,-0.519) -- (0.663,-0.504) -- (0.750,-0.487) -- (0.833,-0.467) -- (0.913,-0.446) -- (0.989,-0.422) -- (1.061,-0.397) -- (1.128,-0.370) -- (1.190,-0.342) -- (1.247,-0.312) -- (1.299,-0.281) -- (1.345,-0.249) -- (1.386,-0.215) -- (1.420,-0.181) -- (1.449,-0.145) -- (1.471,-0.110) -- (1.487,-0.073) -- (1.497,-0.037) -- (1.500,-0.000);
\draw[thick] (-1.500,0.000) -- (-1.500,2.967);
\draw[thick] (1.500,0.000) -- (1.500,2.967);
\draw[fill=teal!26,draw=black,thick] (1.500,2.967) -- (1.487,3.040) -- (1.449,3.112) -- (1.386,3.182) -- (1.299,3.248) -- (1.190,3.309) -- (1.061,3.364) -- (0.913,3.413) -- (0.750,3.454) -- (0.574,3.486) -- (0.388,3.510) -- (0.196,3.524) -- (0.000,3.529) -- (-0.196,3.524) -- (-0.388,3.510) -- (-0.574,3.486) -- (-0.750,3.454) -- (-0.913,3.413) -- (-1.061,3.364) -- (-1.190,3.309) -- (-1.299,3.248) -- (-1.386,3.182) -- (-1.449,3.112) -- (-1.487,3.040) -- (-1.500,2.967) -- (-1.487,2.894) -- (-1.449,2.822) -- (-1.386,2.752) -- (-1.299,2.686) -- (-1.190,2.625) -- (-1.061,2.570) -- (-0.913,2.521) -- (-0.750,2.480) -- (-0.574,2.448) -- (-0.388,2.424) -- (-0.196,2.410) -- (-0.000,2.405) -- (0.196,2.410) -- (0.388,2.424) -- (0.574,2.448) -- (0.750,2.480) -- (0.913,2.521) -- (1.061,2.570) -- (1.190,2.625) -- (1.299,2.686) -- (1.386,2.752) -- (1.449,2.822) -- (1.487,2.894) -- (1.500,2.967) -- cycle;
\draw[dashed,draw=accent,thick] (0.000,0.000) -- (1.500,0.000);
\fill (0.000,0.000) circle (1.3pt);
\node[font=\small,inner sep=2pt,above,text=accent] at (0.750,0.000) {$3$\,cm};
\node[font=\small,inner sep=2pt,right] at (1.500,1.483) {$10$\,cm};
\end{tikzpicture}\end{minipage}\par\recorda{Cilindro: $A_L=2\pi rh$, \ $A_T=2\pi rh+2\pi r^2$. \ Cono: $g^2=h^2+r^2$, \ $A_L=\pi rg$, \ $A_T=\pi rg+\pi r^2$. Usa la tecla $\pi$ y redondea a las centésimas.}

\ex{1} \begin{minipage}[t]{0.61\linewidth}Calcula la generatriz, el área lateral y el área total del cono de la figura.\end{minipage}\hfill\begin{minipage}[t]{0.25\linewidth}\vspace{0pt}\centering \begin{tikzpicture}[x=0.5cm,y=0.5cm,baseline=(current bounding box.center),line join=round]
\path[fill=teal!14,draw=none] (0.000,3.152) -- (-1.571,0.114) -- (-1.589,0.070) -- (-1.598,0.027) -- (-1.599,-0.017) -- (-1.592,-0.061) -- (-1.575,-0.105) -- (-1.551,-0.148) -- (-1.517,-0.190) -- (-1.476,-0.231) -- (-1.427,-0.271) -- (-1.370,-0.310) -- (-1.306,-0.346) -- (-1.234,-0.381) -- (-1.156,-0.414) -- (-1.072,-0.445) -- (-0.982,-0.473) -- (-0.887,-0.499) -- (-0.787,-0.522) -- (-0.682,-0.542) -- (-0.574,-0.559) -- (-0.463,-0.574) -- (-0.350,-0.585) -- (-0.234,-0.593) -- (-0.117,-0.598) -- (-0.000,-0.599) -- (0.117,-0.598) -- (0.234,-0.593) -- (0.350,-0.585) -- (0.463,-0.574) -- (0.574,-0.559) -- (0.682,-0.542) -- (0.787,-0.522) -- (0.887,-0.499) -- (0.982,-0.473) -- (1.072,-0.445) -- (1.156,-0.414) -- (1.234,-0.381) -- (1.306,-0.346) -- (1.370,-0.310) -- (1.427,-0.271) -- (1.476,-0.231) -- (1.517,-0.190) -- (1.551,-0.148) -- (1.575,-0.105) -- (1.592,-0.061) -- (1.599,-0.017) -- (1.598,0.027) -- (1.589,0.070) -- (1.571,0.114) -- cycle;
\draw[dashed,draw=muted] (1.571,0.114) -- (1.551,0.148) -- (1.526,0.181) -- (1.495,0.213) -- (1.460,0.245) -- (1.420,0.276) -- (1.375,0.306) -- (1.326,0.335) -- (1.273,0.363) -- (1.215,0.390) -- (1.153,0.416) -- (1.087,0.440) -- (1.018,0.462) -- (0.946,0.484) -- (0.870,0.503) -- (0.791,0.521) -- (0.710,0.537) -- (0.627,0.552) -- (0.541,0.564) -- (0.454,0.575) -- (0.365,0.584) -- (0.275,0.590) -- (0.184,0.595) -- (0.092,0.598) -- (0.000,0.599) -- (-0.092,0.598) -- (-0.184,0.595) -- (-0.275,0.590) -- (-0.365,0.584) -- (-0.454,0.575) -- (-0.541,0.564) -- (-0.627,0.552) -- (-0.710,0.537) -- (-0.791,0.521) -- (-0.870,0.503) -- (-0.946,0.484) -- (-1.018,0.462) -- (-1.087,0.440) -- (-1.153,0.416) -- (-1.215,0.390) -- (-1.273,0.363) -- (-1.326,0.335) -- (-1.375,0.306) -- (-1.420,0.276) -- (-1.460,0.245) -- (-1.495,0.213) -- (-1.526,0.181) -- (-1.551,0.148) -- (-1.571,0.114);
\draw[thick] (-1.571,0.114) -- (-1.589,0.070) -- (-1.598,0.027) -- (-1.599,-0.017) -- (-1.592,-0.061) -- (-1.575,-0.105) -- (-1.551,-0.148) -- (-1.517,-0.190) -- (-1.476,-0.231) -- (-1.427,-0.271) -- (-1.370,-0.310) -- (-1.306,-0.346) -- (-1.234,-0.381) -- (-1.156,-0.414) -- (-1.072,-0.445) -- (-0.982,-0.473) -- (-0.887,-0.499) -- (-0.787,-0.522) -- (-0.682,-0.542) -- (-0.574,-0.559) -- (-0.463,-0.574) -- (-0.350,-0.585) -- (-0.234,-0.593) -- (-0.117,-0.598) -- (-0.000,-0.599) -- (0.117,-0.598) -- (0.234,-0.593) -- (0.350,-0.585) -- (0.463,-0.574) -- (0.574,-0.559) -- (0.682,-0.542) -- (0.787,-0.522) -- (0.887,-0.499) -- (0.982,-0.473) -- (1.072,-0.445) -- (1.156,-0.414) -- (1.234,-0.381) -- (1.306,-0.346) -- (1.370,-0.310) -- (1.427,-0.271) -- (1.476,-0.231) -- (1.517,-0.190) -- (1.551,-0.148) -- (1.575,-0.105) -- (1.592,-0.061) -- (1.599,-0.017) -- (1.598,0.027) -- (1.589,0.070) -- (1.571,0.114);
\draw[thick] (0.000,3.152) -- (-1.571,0.114);
\draw[thick] (0.000,3.152) -- (1.571,0.114);
\draw[dashed,draw=accent,thick] (0.000,0.000) -- (1.600,0.000);
\draw[dashed,draw=accent,thick] (0.000,0.000) -- (0.000,3.152);
\fill (0.000,0.000) circle (1.3pt);
\node[font=\small,inner sep=2pt,above,text=accent] at (0.800,0.000) {$5$\,cm};
\node[font=\small,inner sep=2pt,left,text=accent] at (0.000,1.576) {$12$\,cm};
\node[font=\small,inner sep=2pt,right] at (0.800,1.576) {$g$};
\end{tikzpicture}\end{minipage}\par

\ex{1} Un cono tiene 6 cm de radio y 10 cm de generatriz. ¿Qué altura tiene? Calcula su área total.

\ex{2} Una lata de refresco tiene 8 cm de diámetro y 12 cm de altura. ¿Cuánta chapa hace falta para fabricarla?

\ex{2} Queremos hacer un gorro de payaso en forma de cono de 9 cm de radio y 12 cm de altura. ¿Cuánto papel necesitamos?

\ex{2} Un cilindro tiene un área lateral de $120\pi$ cm² y 10 cm de altura. ¿Qué radio tiene? Calcula su área total.

\ex{3} Un cono tiene 4 cm de radio y 12 cm de generatriz. ¿Qué ángulo tiene el sector circular de su desarrollo?

\ex{3} Hay que pintar la superficie lateral de 12 columnas cilíndricas de 0,4 m de radio y 5 m de altura. Si un litro de pintura cubre 4 m², ¿cuántos litros hacen falta?


% =====================================================
\seccio{5}{La esfera. La Tierra}

\ex{1} \begin{minipage}[t]{0.64\linewidth}Calcula el área de:\begin{tasks}(1)
  \task una esfera de 4 cm de radio;
  \task una pelota de 10 cm de diámetro;
  \task una bola de 0,5 m de radio.
\end{tasks}\end{minipage}\hfill\begin{minipage}[t]{0.22\linewidth}\vspace{0pt}\centering \begin{tikzpicture}[x=0.5cm,y=0.5cm,baseline=(current bounding box.center),line join=round]
\draw[fill=teal!14,draw=black,thick] (0.000,0.000) circle[radius=1.6];
\draw[dashed,draw=muted] (1.600,0.000) -- (1.597,0.039) -- (1.586,0.078) -- (1.569,0.117) -- (1.545,0.155) -- (1.515,0.193) -- (1.478,0.229) -- (1.435,0.265) -- (1.386,0.300) -- (1.330,0.333) -- (1.269,0.365) -- (1.203,0.395) -- (1.131,0.424) -- (1.055,0.451) -- (0.974,0.476) -- (0.889,0.498) -- (0.800,0.519) -- (0.708,0.538) -- (0.612,0.554) -- (0.514,0.568) -- (0.414,0.579) -- (0.312,0.588) -- (0.209,0.594) -- (0.105,0.598) -- (0.000,0.599) -- (-0.105,0.598) -- (-0.209,0.594) -- (-0.312,0.588) -- (-0.414,0.579) -- (-0.514,0.568) -- (-0.612,0.554) -- (-0.708,0.538) -- (-0.800,0.519) -- (-0.889,0.498) -- (-0.974,0.476) -- (-1.055,0.451) -- (-1.131,0.424) -- (-1.203,0.395) -- (-1.269,0.365) -- (-1.330,0.333) -- (-1.386,0.300) -- (-1.435,0.265) -- (-1.478,0.229) -- (-1.515,0.193) -- (-1.545,0.155) -- (-1.569,0.117) -- (-1.586,0.078) -- (-1.597,0.039) -- (-1.600,0.000);
\draw[thick] (-1.600,0.000) -- (-1.597,-0.039) -- (-1.586,-0.078) -- (-1.569,-0.117) -- (-1.545,-0.155) -- (-1.515,-0.193) -- (-1.478,-0.229) -- (-1.435,-0.265) -- (-1.386,-0.300) -- (-1.330,-0.333) -- (-1.269,-0.365) -- (-1.203,-0.395) -- (-1.131,-0.424) -- (-1.055,-0.451) -- (-0.974,-0.476) -- (-0.889,-0.498) -- (-0.800,-0.519) -- (-0.708,-0.538) -- (-0.612,-0.554) -- (-0.514,-0.568) -- (-0.414,-0.579) -- (-0.312,-0.588) -- (-0.209,-0.594) -- (-0.105,-0.598) -- (-0.000,-0.599) -- (0.105,-0.598) -- (0.209,-0.594) -- (0.312,-0.588) -- (0.414,-0.579) -- (0.514,-0.568) -- (0.612,-0.554) -- (0.708,-0.538) -- (0.800,-0.519) -- (0.889,-0.498) -- (0.974,-0.476) -- (1.055,-0.451) -- (1.131,-0.424) -- (1.203,-0.395) -- (1.269,-0.365) -- (1.330,-0.333) -- (1.386,-0.300) -- (1.435,-0.265) -- (1.478,-0.229) -- (1.515,-0.193) -- (1.545,-0.155) -- (1.569,-0.117) -- (1.586,-0.078) -- (1.597,-0.039) -- (1.600,-0.000);
\draw[dashed,draw=accent,thick] (0.000,0.000) -- (1.600,0.000);
\fill (0.000,0.000) circle (1.3pt);
\node[font=\small,inner sep=2pt,above,text=accent] at (0.800,0.000) {$r$};
\end{tikzpicture}\end{minipage}\par\recorda{Área de la esfera: $A=4\pi r^2$ (cuatro círculos máximos). \ La \textbf{latitud} va de $0^\circ$ a $90^\circ$ N o S (desde el ecuador) y la \textbf{longitud}, de $0^\circ$ a $180^\circ$ E u O (desde el meridiano de Greenwich).}

\ex{2} Una esfera tiene un área de $144\pi$ cm². ¿Qué radio tiene?

\ex{2} Una cúpula tiene forma de semiesfera de 5 m de radio. ¿Qué superficie hay que impermeabilizar por fuera?

\ex{2} Dos ciudades están sobre el mismo meridiano, una a $40^\circ$ N y la otra a $10^\circ$ S. ¿Cuántos grados de latitud las separan? Si un meridiano entero (la vuelta a la Tierra) mide 40\,000 km, ¿qué distancia hay entre ellas sobre el meridiano?

\ex{2} Cuando en Greenwich son las 12 horas solares, ¿qué hora solar es en un lugar de longitud $45^\circ$ E? ¿Y en uno de longitud $30^\circ$ O?

\ex{3} El radio de la Tierra es 6371 km y el de la Luna, 1737 km. ¿Cuántas veces mayor es la superficie de la Tierra que la de la Luna?


% =====================================================
\seccio{6}{Volumen. Prismas y cilindros}

\ex{1} Expresa en la unidad indicada:\begin{tasks}(3)
  \task 3,2 m³ en L
  \task 450 cm³ en L
  \task 2,5 L en cm³
  \task 0,75 dm³ en mL
  \task 12\,000 L en m³
  \task 5 hL en m³
\end{tasks}\recorda{Cada unidad de volumen es 1000 veces la siguiente: $1\text{ m}^3=1000\text{ dm}^3$, \ $1\text{ dm}^3=1000\text{ cm}^3$. \ $1\text{ dm}^3=1\text{ L}$, \ $1\text{ cm}^3=1\text{ mL}$. \ Prisma y cilindro: $V=A_{\text{base}}\cdot h$; \ $V_{\text{cilindro}}=\pi r^2h$.}

\ex{1} Calcula el volumen de:\begin{tasks}(1)
  \task un cubo de 6 cm de arista;
  \task un ortoedro de 8 cm × 5 cm × 3 cm;
  \task un cilindro de 4 cm de radio y 10 cm de altura.
\end{tasks}

\ex{1} La base de un prisma es un triángulo rectángulo de catetos 5 cm y 12 cm. Si el prisma mide 20 cm de altura, ¿qué volumen tiene?

\ex{2} Calcula el volumen de un prisma hexagonal regular de 2 cm de arista de la base y 5 cm de altura.

\ex{2} Un acuario mide 60 cm × 30 cm de base y 40 cm de altura. ¿Cuántos litros caben? ¿Cuántos litros hay si el agua llega a $\frac34$ de la altura?

\ex{2} Un depósito cilíndrico tiene 1,5 m de radio y 4 m de altura. ¿Cuántos litros de agua caben?

\ex{2} Una piscina mide 25 m de largo y 10 m de ancho. El fondo es inclinado: la profundidad va de 1 m a 3 m. ¿Qué volumen de agua cabe? (Es un prisma con una base en forma de trapecio.)

\ex{3} Un vaso cilíndrico mide 4 cm de radio. Si echamos medio litro de agua, ¿hasta qué altura llega?

\ex{3} Un ladrillo mide 24 cm × 11,5 cm × 5 cm y tiene tres agujeros cilíndricos de 1,5 cm de radio que lo atraviesan (en la dirección de los 5 cm). ¿Qué volumen de arcilla tiene?


% =====================================================
\seccio{7}{Volumen de pirámides, conos y esferas}

\ex{1} Calcula el volumen de:\begin{tasks}(1)
  \task una pirámide cuadrangular de 6 cm de arista de la base y 10 cm de altura;
  \task un cono de 3 cm de radio y 7 cm de altura;
  \task una esfera de 6 cm de radio.
\end{tasks}\recorda{$V_{\text{pirámide}}=\frac{A_{\text{base}}\cdot h}{3}$; \ $V_{\text{cono}}=\frac{\pi r^2h}{3}$; \ $V_{\text{esfera}}=\frac43\pi r^3$.}

\ex{1} ¿Qué volumen tiene una bola de 20 cm de diámetro? Exprésalo en cm³ y en litros.

\ex{2} Una pirámide cuadrangular regular tiene 10 cm de arista de la base y 13 cm de apotema. Calcula su altura y su volumen.

\ex{2} Un cono tiene 6 cm de radio y 10 cm de generatriz. Calcula su volumen.

\ex{2} ¿Cuántos litros de aire hay dentro de una pelota de baloncesto de 11 cm de radio?

\ex{2} Calcula el volumen de una pirámide hexagonal regular de 4 cm de arista de la base y 9 cm de altura.

\ex{3} Un bol tiene forma de semiesfera de 12 cm de radio interior. ¿Cuántos litros caben?

\ex{3} Una esfera y un cono tienen el mismo radio, 6 cm, y el mismo volumen. ¿Qué altura tiene el cono?


% =====================================================
\seccio{8}{Cuerpos compuestos y problemas}

\ex{1} Una pieza está formada por un cubo de 6 cm de arista y, encima, una pirámide de la misma base y 4 cm de altura. Calcula su volumen.\recorda{Descompón el cuerpo en cuerpos simples. El volumen es la suma (o la diferencia, si hay agujeros) de los volúmenes. Para el área, suma solo las superficies que quedan por fuera.}

\ex{2} \begin{minipage}[t]{0.61\linewidth}Un silo tiene la forma de la figura: un cilindro de 3 m de radio y 10 m de altura con un tejado cónico de 4 m de altura. ¿Qué volumen tiene?\end{minipage}\hfill\begin{minipage}[t]{0.25\linewidth}\vspace{0pt}\centering \begin{tikzpicture}[x=0.42cm,y=0.42cm,baseline=(current bounding box.center),line join=round]
\path[fill=teal!14,draw=none] (-1.500,3.152) -- (-1.500,0.000) -- (-1.497,-0.037) -- (-1.487,-0.073) -- (-1.471,-0.110) -- (-1.449,-0.145) -- (-1.420,-0.181) -- (-1.386,-0.215) -- (-1.345,-0.249) -- (-1.299,-0.281) -- (-1.247,-0.312) -- (-1.190,-0.342) -- (-1.128,-0.370) -- (-1.061,-0.397) -- (-0.989,-0.422) -- (-0.913,-0.446) -- (-0.833,-0.467) -- (-0.750,-0.487) -- (-0.663,-0.504) -- (-0.574,-0.519) -- (-0.482,-0.532) -- (-0.388,-0.543) -- (-0.293,-0.551) -- (-0.196,-0.557) -- (-0.098,-0.561) -- (-0.000,-0.562) -- (0.098,-0.561) -- (0.196,-0.557) -- (0.293,-0.551) -- (0.388,-0.543) -- (0.482,-0.532) -- (0.574,-0.519) -- (0.663,-0.504) -- (0.750,-0.487) -- (0.833,-0.467) -- (0.913,-0.446) -- (0.989,-0.422) -- (1.061,-0.397) -- (1.128,-0.370) -- (1.190,-0.342) -- (1.247,-0.312) -- (1.299,-0.281) -- (1.345,-0.249) -- (1.386,-0.215) -- (1.420,-0.181) -- (1.449,-0.145) -- (1.471,-0.110) -- (1.487,-0.073) -- (1.497,-0.037) -- (1.500,-0.000) -- (1.500,3.152) -- (1.500,3.152) -- (1.497,3.116) -- (1.487,3.079) -- (1.471,3.043) -- (1.449,3.007) -- (1.420,2.972) -- (1.386,2.937) -- (1.345,2.904) -- (1.299,2.871) -- (1.247,2.840) -- (1.190,2.810) -- (1.128,2.782) -- (1.061,2.755) -- (0.989,2.730) -- (0.913,2.707) -- (0.833,2.685) -- (0.750,2.666) -- (0.663,2.648) -- (0.574,2.633) -- (0.482,2.620) -- (0.388,2.610) -- (0.293,2.601) -- (0.196,2.595) -- (0.098,2.592) -- (-0.000,2.591) -- (-0.098,2.592) -- (-0.196,2.595) -- (-0.293,2.601) -- (-0.388,2.610) -- (-0.482,2.620) -- (-0.574,2.633) -- (-0.663,2.648) -- (-0.750,2.666) -- (-0.833,2.685) -- (-0.913,2.707) -- (-0.989,2.730) -- (-1.061,2.755) -- (-1.128,2.782) -- (-1.190,2.810) -- (-1.247,2.840) -- (-1.299,2.871) -- (-1.345,2.904) -- (-1.386,2.937) -- (-1.420,2.972) -- (-1.449,3.007) -- (-1.471,3.043) -- (-1.487,3.079) -- (-1.497,3.116) -- (-1.500,3.152) -- cycle;
\draw[dashed,draw=muted] (1.500,0.000) -- (1.497,0.037) -- (1.487,0.073) -- (1.471,0.110) -- (1.449,0.145) -- (1.420,0.181) -- (1.386,0.215) -- (1.345,0.249) -- (1.299,0.281) -- (1.247,0.312) -- (1.190,0.342) -- (1.128,0.370) -- (1.061,0.397) -- (0.989,0.422) -- (0.913,0.446) -- (0.833,0.467) -- (0.750,0.487) -- (0.663,0.504) -- (0.574,0.519) -- (0.482,0.532) -- (0.388,0.543) -- (0.293,0.551) -- (0.196,0.557) -- (0.098,0.561) -- (0.000,0.562) -- (-0.098,0.561) -- (-0.196,0.557) -- (-0.293,0.551) -- (-0.388,0.543) -- (-0.482,0.532) -- (-0.574,0.519) -- (-0.663,0.504) -- (-0.750,0.487) -- (-0.833,0.467) -- (-0.913,0.446) -- (-0.989,0.422) -- (-1.061,0.397) -- (-1.128,0.370) -- (-1.190,0.342) -- (-1.247,0.312) -- (-1.299,0.281) -- (-1.345,0.249) -- (-1.386,0.215) -- (-1.420,0.181) -- (-1.449,0.145) -- (-1.471,0.110) -- (-1.487,0.073) -- (-1.497,0.037) -- (-1.500,0.000);
\draw[thick] (-1.500,0.000) -- (-1.497,-0.037) -- (-1.487,-0.073) -- (-1.471,-0.110) -- (-1.449,-0.145) -- (-1.420,-0.181) -- (-1.386,-0.215) -- (-1.345,-0.249) -- (-1.299,-0.281) -- (-1.247,-0.312) -- (-1.190,-0.342) -- (-1.128,-0.370) -- (-1.061,-0.397) -- (-0.989,-0.422) -- (-0.913,-0.446) -- (-0.833,-0.467) -- (-0.750,-0.487) -- (-0.663,-0.504) -- (-0.574,-0.519) -- (-0.482,-0.532) -- (-0.388,-0.543) -- (-0.293,-0.551) -- (-0.196,-0.557) -- (-0.098,-0.561) -- (-0.000,-0.562) -- (0.098,-0.561) -- (0.196,-0.557) -- (0.293,-0.551) -- (0.388,-0.543) -- (0.482,-0.532) -- (0.574,-0.519) -- (0.663,-0.504) -- (0.750,-0.487) -- (0.833,-0.467) -- (0.913,-0.446) -- (0.989,-0.422) -- (1.061,-0.397) -- (1.128,-0.370) -- (1.190,-0.342) -- (1.247,-0.312) -- (1.299,-0.281) -- (1.345,-0.249) -- (1.386,-0.215) -- (1.420,-0.181) -- (1.449,-0.145) -- (1.471,-0.110) -- (1.487,-0.073) -- (1.497,-0.037) -- (1.500,-0.000);
\draw[thick] (-1.500,0.000) -- (-1.500,3.152);
\draw[thick] (1.500,0.000) -- (1.500,3.152);
\path[fill=teal!14,draw=none] (0.000,4.450) -- (-1.352,3.396) -- (-1.402,3.352) -- (-1.442,3.307) -- (-1.471,3.262) -- (-1.491,3.215) -- (-1.499,3.168) -- (-1.498,3.120) -- (-1.485,3.073) -- (-1.462,3.027) -- (-1.429,2.981) -- (-1.385,2.937) -- (-1.332,2.894) -- (-1.270,2.853) -- (-1.198,2.814) -- (-1.118,2.778) -- (-1.030,2.744) -- (-0.935,2.713) -- (-0.833,2.685) -- (-0.725,2.661) -- (-0.612,2.639) -- (-0.495,2.622) -- (-0.374,2.608) -- (-0.251,2.598) -- (-0.126,2.593) -- (-0.000,2.591) -- (0.126,2.593) -- (0.251,2.598) -- (0.374,2.608) -- (0.495,2.622) -- (0.612,2.639) -- (0.725,2.661) -- (0.833,2.685) -- (0.935,2.713) -- (1.030,2.744) -- (1.118,2.778) -- (1.198,2.814) -- (1.270,2.853) -- (1.332,2.894) -- (1.385,2.937) -- (1.429,2.981) -- (1.462,3.027) -- (1.485,3.073) -- (1.498,3.120) -- (1.499,3.168) -- (1.491,3.215) -- (1.471,3.262) -- (1.442,3.307) -- (1.402,3.352) -- (1.352,3.396) -- cycle;
\draw[dashed,draw=muted] (1.352,3.396) -- (1.320,3.419) -- (1.286,3.442) -- (1.248,3.464) -- (1.208,3.486) -- (1.165,3.506) -- (1.119,3.527) -- (1.071,3.546) -- (1.021,3.564) -- (0.969,3.581) -- (0.914,3.598) -- (0.857,3.614) -- (0.799,3.628) -- (0.738,3.642) -- (0.677,3.654) -- (0.613,3.665) -- (0.549,3.675) -- (0.483,3.684) -- (0.416,3.692) -- (0.348,3.699) -- (0.279,3.705) -- (0.210,3.709) -- (0.140,3.712) -- (0.070,3.714) -- (0.000,3.714) -- (-0.070,3.714) -- (-0.140,3.712) -- (-0.210,3.709) -- (-0.279,3.705) -- (-0.348,3.699) -- (-0.416,3.692) -- (-0.483,3.684) -- (-0.549,3.675) -- (-0.613,3.665) -- (-0.677,3.654) -- (-0.738,3.642) -- (-0.799,3.628) -- (-0.857,3.614) -- (-0.914,3.598) -- (-0.969,3.581) -- (-1.021,3.564) -- (-1.071,3.546) -- (-1.119,3.527) -- (-1.165,3.506) -- (-1.208,3.486) -- (-1.248,3.464) -- (-1.286,3.442) -- (-1.320,3.419) -- (-1.352,3.396);
\draw[thick] (-1.352,3.396) -- (-1.402,3.352) -- (-1.442,3.307) -- (-1.471,3.262) -- (-1.491,3.215) -- (-1.499,3.168) -- (-1.498,3.120) -- (-1.485,3.073) -- (-1.462,3.027) -- (-1.429,2.981) -- (-1.385,2.937) -- (-1.332,2.894) -- (-1.270,2.853) -- (-1.198,2.814) -- (-1.118,2.778) -- (-1.030,2.744) -- (-0.935,2.713) -- (-0.833,2.685) -- (-0.725,2.661) -- (-0.612,2.639) -- (-0.495,2.622) -- (-0.374,2.608) -- (-0.251,2.598) -- (-0.126,2.593) -- (-0.000,2.591) -- (0.126,2.593) -- (0.251,2.598) -- (0.374,2.608) -- (0.495,2.622) -- (0.612,2.639) -- (0.725,2.661) -- (0.833,2.685) -- (0.935,2.713) -- (1.030,2.744) -- (1.118,2.778) -- (1.198,2.814) -- (1.270,2.853) -- (1.332,2.894) -- (1.385,2.937) -- (1.429,2.981) -- (1.462,3.027) -- (1.485,3.073) -- (1.498,3.120) -- (1.499,3.168) -- (1.491,3.215) -- (1.471,3.262) -- (1.442,3.307) -- (1.402,3.352) -- (1.352,3.396);
\draw[thick] (0.000,4.450) -- (-1.352,3.396);
\draw[thick] (0.000,4.450) -- (1.352,3.396);
\draw[dashed,draw=accent,thick] (0.000,0.000) -- (1.500,0.000);
\node[font=\small,inner sep=2pt,below,text=accent] at (0.750,0.000) {$3$\,m};
\node[font=\small,inner sep=2pt,right] at (1.500,1.576) {$10$\,m};
\draw[dashed,draw=accent,thick] (0.000,3.152) -- (0.000,4.450);
\node[font=\small,inner sep=2pt,left,text=accent] at (-0.800,3.801) {$4$\,m};
\end{tikzpicture}\end{minipage}\par

\ex{2} \begin{minipage}[t]{0.66\linewidth}Una cápsula es un cilindro de 2 cm de radio y 6 cm de largo con una semiesfera en cada extremo. Calcula su volumen y su área.\end{minipage}\hfill\begin{minipage}[t]{0.2\linewidth}\vspace{0pt}\centering \begin{tikzpicture}[x=0.5cm,y=0.5cm,baseline=(current bounding box.center),line join=round]
\path[fill=teal!14,draw=none] (1.000,2.782) -- (0.998,2.847) -- (0.991,2.912) -- (0.981,2.977) -- (0.966,3.040) -- (0.947,3.103) -- (0.924,3.164) -- (0.897,3.224) -- (0.866,3.282) -- (0.831,3.337) -- (0.793,3.390) -- (0.752,3.441) -- (0.707,3.489) -- (0.659,3.533) -- (0.609,3.575) -- (0.556,3.613) -- (0.500,3.648) -- (0.442,3.678) -- (0.383,3.705) -- (0.321,3.728) -- (0.259,3.747) -- (0.195,3.762) -- (0.131,3.773) -- (0.065,3.779) -- (0.000,3.782) -- (-0.065,3.779) -- (-0.131,3.773) -- (-0.195,3.762) -- (-0.259,3.747) -- (-0.321,3.728) -- (-0.383,3.705) -- (-0.442,3.678) -- (-0.500,3.648) -- (-0.556,3.613) -- (-0.609,3.575) -- (-0.659,3.533) -- (-0.707,3.489) -- (-0.752,3.441) -- (-0.793,3.390) -- (-0.831,3.337) -- (-0.866,3.282) -- (-0.897,3.224) -- (-0.924,3.164) -- (-0.947,3.103) -- (-0.966,3.040) -- (-0.981,2.977) -- (-0.991,2.912) -- (-0.998,2.847) -- (-1.000,2.782) -- (-1.000,-0.000) -- (-0.998,-0.065) -- (-0.991,-0.131) -- (-0.981,-0.195) -- (-0.966,-0.259) -- (-0.947,-0.321) -- (-0.924,-0.383) -- (-0.897,-0.442) -- (-0.866,-0.500) -- (-0.831,-0.556) -- (-0.793,-0.609) -- (-0.752,-0.659) -- (-0.707,-0.707) -- (-0.659,-0.752) -- (-0.609,-0.793) -- (-0.556,-0.831) -- (-0.500,-0.866) -- (-0.442,-0.897) -- (-0.383,-0.924) -- (-0.321,-0.947) -- (-0.259,-0.966) -- (-0.195,-0.981) -- (-0.131,-0.991) -- (-0.065,-0.998) -- (-0.000,-1.000) -- (0.065,-0.998) -- (0.131,-0.991) -- (0.195,-0.981) -- (0.259,-0.966) -- (0.321,-0.947) -- (0.383,-0.924) -- (0.442,-0.897) -- (0.500,-0.866) -- (0.556,-0.831) -- (0.609,-0.793) -- (0.659,-0.752) -- (0.707,-0.707) -- (0.752,-0.659) -- (0.793,-0.609) -- (0.831,-0.556) -- (0.866,-0.500) -- (0.897,-0.442) -- (0.924,-0.383) -- (0.947,-0.321) -- (0.966,-0.259) -- (0.981,-0.195) -- (0.991,-0.131) -- (0.998,-0.065) -- (1.000,-0.000) -- cycle;
\draw[thick] (-1.000,0.000) -- (-0.998,-0.025) -- (-0.991,-0.049) -- (-0.981,-0.073) -- (-0.966,-0.097) -- (-0.947,-0.120) -- (-0.924,-0.143) -- (-0.897,-0.166) -- (-0.866,-0.187) -- (-0.831,-0.208) -- (-0.793,-0.228) -- (-0.752,-0.247) -- (-0.707,-0.265) -- (-0.659,-0.282) -- (-0.609,-0.297) -- (-0.556,-0.311) -- (-0.500,-0.324) -- (-0.442,-0.336) -- (-0.383,-0.346) -- (-0.321,-0.355) -- (-0.259,-0.362) -- (-0.195,-0.367) -- (-0.131,-0.371) -- (-0.065,-0.374) -- (-0.000,-0.375) -- (0.065,-0.374) -- (0.131,-0.371) -- (0.195,-0.367) -- (0.259,-0.362) -- (0.321,-0.355) -- (0.383,-0.346) -- (0.442,-0.336) -- (0.500,-0.324) -- (0.556,-0.311) -- (0.609,-0.297) -- (0.659,-0.282) -- (0.707,-0.265) -- (0.752,-0.247) -- (0.793,-0.228) -- (0.831,-0.208) -- (0.866,-0.187) -- (0.897,-0.166) -- (0.924,-0.143) -- (0.947,-0.120) -- (0.966,-0.097) -- (0.981,-0.073) -- (0.991,-0.049) -- (0.998,-0.025) -- (1.000,-0.000);
\draw[thick] (-1.000,2.782) -- (-0.998,2.757) -- (-0.991,2.733) -- (-0.981,2.708) -- (-0.966,2.685) -- (-0.947,2.661) -- (-0.924,2.638) -- (-0.897,2.616) -- (-0.866,2.594) -- (-0.831,2.573) -- (-0.793,2.554) -- (-0.752,2.535) -- (-0.707,2.517) -- (-0.659,2.500) -- (-0.609,2.484) -- (-0.556,2.470) -- (-0.500,2.457) -- (-0.442,2.446) -- (-0.383,2.435) -- (-0.321,2.427) -- (-0.259,2.420) -- (-0.195,2.414) -- (-0.131,2.410) -- (-0.065,2.408) -- (-0.000,2.407) -- (0.065,2.408) -- (0.131,2.410) -- (0.195,2.414) -- (0.259,2.420) -- (0.321,2.427) -- (0.383,2.435) -- (0.442,2.446) -- (0.500,2.457) -- (0.556,2.470) -- (0.609,2.484) -- (0.659,2.500) -- (0.707,2.517) -- (0.752,2.535) -- (0.793,2.554) -- (0.831,2.573) -- (0.866,2.594) -- (0.897,2.616) -- (0.924,2.638) -- (0.947,2.661) -- (0.966,2.685) -- (0.981,2.708) -- (0.991,2.733) -- (0.998,2.757) -- (1.000,2.782);
\draw[thick] (1.000,2.782) -- (0.998,2.847) -- (0.991,2.912) -- (0.981,2.977) -- (0.966,3.040) -- (0.947,3.103) -- (0.924,3.164) -- (0.897,3.224) -- (0.866,3.282) -- (0.831,3.337) -- (0.793,3.390) -- (0.752,3.441) -- (0.707,3.489) -- (0.659,3.533) -- (0.609,3.575) -- (0.556,3.613) -- (0.500,3.648) -- (0.442,3.678) -- (0.383,3.705) -- (0.321,3.728) -- (0.259,3.747) -- (0.195,3.762) -- (0.131,3.773) -- (0.065,3.779) -- (0.000,3.782) -- (-0.065,3.779) -- (-0.131,3.773) -- (-0.195,3.762) -- (-0.259,3.747) -- (-0.321,3.728) -- (-0.383,3.705) -- (-0.442,3.678) -- (-0.500,3.648) -- (-0.556,3.613) -- (-0.609,3.575) -- (-0.659,3.533) -- (-0.707,3.489) -- (-0.752,3.441) -- (-0.793,3.390) -- (-0.831,3.337) -- (-0.866,3.282) -- (-0.897,3.224) -- (-0.924,3.164) -- (-0.947,3.103) -- (-0.966,3.040) -- (-0.981,2.977) -- (-0.991,2.912) -- (-0.998,2.847) -- (-1.000,2.782) -- (-1.000,-0.000) -- (-0.998,-0.065) -- (-0.991,-0.131) -- (-0.981,-0.195) -- (-0.966,-0.259) -- (-0.947,-0.321) -- (-0.924,-0.383) -- (-0.897,-0.442) -- (-0.866,-0.500) -- (-0.831,-0.556) -- (-0.793,-0.609) -- (-0.752,-0.659) -- (-0.707,-0.707) -- (-0.659,-0.752) -- (-0.609,-0.793) -- (-0.556,-0.831) -- (-0.500,-0.866) -- (-0.442,-0.897) -- (-0.383,-0.924) -- (-0.321,-0.947) -- (-0.259,-0.966) -- (-0.195,-0.981) -- (-0.131,-0.991) -- (-0.065,-0.998) -- (-0.000,-1.000) -- (0.065,-0.998) -- (0.131,-0.991) -- (0.195,-0.981) -- (0.259,-0.966) -- (0.321,-0.947) -- (0.383,-0.924) -- (0.442,-0.897) -- (0.500,-0.866) -- (0.556,-0.831) -- (0.609,-0.793) -- (0.659,-0.752) -- (0.707,-0.707) -- (0.752,-0.659) -- (0.793,-0.609) -- (0.831,-0.556) -- (0.866,-0.500) -- (0.897,-0.442) -- (0.924,-0.383) -- (0.947,-0.321) -- (0.966,-0.259) -- (0.981,-0.195) -- (0.991,-0.131) -- (0.998,-0.065) -- (1.000,-0.000) -- (1.000,2.782);
\node[font=\small,inner sep=2pt,right] at (1.000,1.391) {$6$\,cm};
\draw[dashed,draw=accent,thick] (0.000,0.000) -- (1.000,0.000);
\node[font=\small,inner sep=2pt,below,text=accent] at (0.500,-1.000) {$2$\,cm};
\end{tikzpicture}\end{minipage}\par

\ex{2} Un helado es un cucurucho cónico de 2,5 cm de radio y 12 cm de altura, lleno, con una semiesfera de helado encima. ¿Qué volumen de helado hay?

\ex{2} Un tubo de 50 cm de largo tiene 6 cm de radio exterior y 5 cm de radio interior. ¿Qué volumen de material tiene?

\ex{2} Un cubo de 10 cm de arista tiene un agujero cilíndrico de 3 cm de radio que lo atraviesa de lado a lado. ¿Qué volumen tiene?

\ex{3} Un cubo tiene forma de tronco de cono: el fondo mide 12 cm de radio, la boca 18 cm y la altura es de 24 cm. ¿Cuántos litros caben? \emph{(Pista: es un cono grande al que le quitamos un cono pequeño; calcula la altura del pequeño con el teorema de Tales.)}

\ex{3} Dejamos caer una bola de acero de 3 cm de radio dentro de un vaso cilíndrico de 6 cm de radio que contiene agua (y la bola queda cubierta). ¿Cuánto sube el nivel del agua?

\ex{3} Un depósito de agua es un cilindro de 2 m de radio y 5 m de altura con el fondo en forma de semiesfera. ¿Cuántos m³ caben? ¿Cuánto tiempo tardará en llenarse con un grifo que vierte 50 L por minuto?

\vfill
\begin{center}\color{muted}\small\faCheckCircle[regular]\ \ ¡Ánimo! Si te cuesta algún ejercicio, repasa los apuntes interactivos de la unidad.\end{center}

\end{document}
